\chapter{Mental health and older adults}

\section*{Definition and Overview}
Mental health and older adults is a mental health condition that affects cognition, emotional regulation, or behaviour. Mental health disorders are defined by clinically significant disturbance in these domains, associated with distress and/or functional impairment. They are common, frequently co-morbid with physical illness, and are a leading cause of disability worldwide. Most are treatable, and many people recover fully or achieve substantial improvement with appropriate intervention.

\section*{Epidemiology}
Approximately one in five adults experiences a mental health condition in any given year. Lifetime prevalence is higher, with anxiety and depressive disorders the most common. Onset is often in adolescence or early adulthood. Co-morbidity with substance use disorders and chronic physical conditions is high. Stigma, delayed help-seeking, and unequal access to care remain significant barriers to effective treatment.

\section*{Aetiology and Pathophysiology}
Mental health conditions arise from a complex interplay of biological, psychological, and social factors:

\begin{itemize}
    \item \textbf{Genetic predisposition:} family history is a strong risk factor; heritability estimates range from 30\% to 80\% depending on the condition
    \item \textbf{Neurobiological factors:} alterations in neurotransmitter systems (serotonin, dopamine, GABA, noradrenaline), neuroendocrine function, and brain circuitry
    \item \textbf{Psychological factors:} maladaptive cognitions, coping styles, and personality traits
    \item \textbf{Early life adversity:} childhood trauma, abuse, and neglect significantly increase lifetime risk
    \item \textbf{Social determinants:} poverty, social isolation, relationship breakdown, unemployment, and chronic stress
    \item \textbf{Substance use:} alcohol and other drugs can trigger or exacerbate symptoms
\end{itemize}

\section*{Clinical Features}
\begin{itemize}
    \item Persistent low mood, sadness, emptiness, or irritability
    \item Excessive worry, fear, panic, or apprehension
    \item Anhedonia -- loss of interest or pleasure in usual activities
    \item Sleep disturbance (insomnia or hypersomnia)
    \item Appetite or weight changes
    \item Fatigue or loss of energy
    \item Psychomotor agitation or retardation
    \item Difficulty concentrating, indecisiveness, or memory complaints
    \item Feelings of worthlessness or excessive guilt
    \item Recurrent thoughts of death or suicide
    \item Hallucinations or delusions (in psychotic disorders)
    \item Avoidance behaviour and social withdrawal
\end{itemize}

\section*{Differential Diagnosis}
\begin{itemize}
    \item Other psychiatric disorders (e.g.\ depression vs.\ bipolar disorder; anxiety vs.\ thyroid dysfunction)
    \item Substance-induced mood or anxiety disorder
    \item Medical conditions mimicking psychiatric presentation (hypothyroidism, anaemia, vitamin B12 deficiency, cerebrovascular disease)
    \item Medication side effects (e.g.\ corticosteroids, beta-blockers)
    \item Normal grief or adjustment reactions
    \item Neurodevelopmental disorders (ADHD, autism spectrum)
\end{itemize}

\section*{When to Seek Medical Care}
Consult a GP or mental health professional if symptoms persist for more than two weeks, cause significant distress, or impair function at work, school, or in relationships.

\textbf{Seek immediate help if there are thoughts of suicide or self-harm.} Call 000, contact Lifeline on 13 11 14, or go to the nearest hospital emergency department. Acute mental health crises are medical emergencies.

\section*{Diagnosis and Investigations}
\begin{itemize}
    \item \textbf{Clinical interview:} thorough history of presenting complaint, psychiatric review of symptoms, and mental state examination
    \item \textbf{Risk assessment:} evaluation of suicide risk, self-harm, and risk to others
    \item \textbf{Physical examination:} to identify medical causes or co-morbidities
    \item \textbf{Blood tests:} full blood count, thyroid function, B12, folate, and metabolic screen to exclude organic causes
    \item \textbf{Screening tools:} standardised questionnaires (e.g.\ PHQ-9 for depression, GAD-7 for anxiety)
    \item \textbf{Psychometric testing:} where personality or cognitive assessment is needed
\end{itemize}

\section*{Management}
\begin{itemize}
    \item \textbf{Psychological therapies:} cognitive behaviour therapy (CBT), interpersonal therapy, dialectical behaviour therapy (DBT), or other evidence-based modalities
    \item \textbf{Pharmacotherapy:} antidepressants, anxiolytics, mood stabilisers, antipsychotics, or other medicines as indicated, prescribed and monitored by a doctor
    \item \textbf{Lifestyle interventions:} regular exercise, sleep hygiene, reduced alcohol and substance use, and a balanced diet
    \item \textbf{Social support:} family involvement, peer support groups, and community services
    \item \textbf{Care coordination:} GP Mental Health Treatment Plan, referral to a psychologist or psychiatrist
    \item \textbf{Acute and crisis care:} hospital admission or community crisis team involvement for severe presentations
\end{itemize}

\section*{Complications}
\begin{itemize}
    \item Suicide and self-harm -- the most serious complication
    \item Substance use disorders (alcohol, drugs) as self-medication
    \item Social and relationship breakdown, isolation
    \item Occupational impairment and unemployment
    \item Physical health decline (poor diet, inactivity, chronic stress)
    \item Stigma and discrimination affecting quality of life
    \item Treatment-resistant symptoms requiring specialist input
\end{itemize}

\section*{Prognosis and Outcomes}
Most people with mental health conditions improve with appropriate treatment. Early intervention is associated with better outcomes. Some conditions (e.g.\ depression, anxiety) may resolve completely; others (e.g.\ schizophrenia, bipolar disorder) typically follow a relapsing-remitting course and require long-term management. Relapse rates are reduced by medication adherence, ongoing therapy, and social support. Approximately 50--70\% of people achieve significant symptom reduction with first-line treatment.

\section*{Prevention and Self-Care}
\begin{itemize}
    \item Build and maintain strong social connections
    \item Develop healthy coping strategies for stress (mindfulness, exercise, relaxation)
    \item Maintain a consistent sleep routine
    \item Limit alcohol and avoid recreational drugs
    \item Seek help early when symptoms first appear
    \item Adhere to treatment plans and attend follow-up appointments
    \item Engage in meaningful activities and set realistic goals
    \item Learn to recognise early warning signs of relapse
\end{itemize}
